\section{Discussion}
\subsection{What the Chain of Evidence Establishes}
The six steps of Table~\ref{tab:roadmap} were designed so that each result raises the question the next one answers, and read together they support one conclusion. The two-step method of Nishimura et al. is exactly reproducible and its central idea is sound, but using the fingerprint estimate as a hard constraint on an incomplete $c$-domain mechanism locks a mechanistic error into the dynamic fit that in-sample statistics cannot detect (Steps~1--2). Relaxing the constraint to a calibrated prior and completing the mechanism with pseudo-processivity turns an in-sample description into an out-of-sample predictor (Step~3), the gain is attributable to exactly those two changes (Step~4), the mechanism recovers the truth where the two-step method is biased and does nothing where it is not needed (Step~5), and the mechanism, not curve-fitting flexibility, is what carries over to extrapolation and to the meaning of the parameters (Step~6). The following subsections take up the individual lessons.

\subsection{What the Reproduction Audit Adds}
Because Nishimura et al. published every data point, their study could be audited to a level that is rarely possible: \NCheckPass{} of \NCheckAll{} data checks and all four fingerprinting fits were reproduced exactly, and the printed parameters of \ReproConsistent{} of the \ReproAll{} dynamic fits recover the printed objective values, which also reveals the enzyme loadings hidden in the product $k_A[E]_T$ (2.0, 5.0 and 15.0\,\%\,w/w) and the actual initial acid:TMP ratio of the Mao case (10:3.3 rather than the nominal 3:1). The audit also shows two methodological problems that are independent of our new model. First, the full dynamic models are not identifiable: the profile of $F_\mathrm{obj}$ along $F_A$ decreases monotonically towards $F_A\to\infty$ (Figure~\ref{fig:landscape}b), the limit in which the rate becomes independent of acid saturation and $k_E$ and the $F$ constants can only be estimated as ratios \cite{raue2009,transtrum2015}; the very large printed confidence intervals of the full models are a symptom of this ridge, and so is the erratic extrapolation of ORIG-full in Step~3. Second, several printed full-model fits are not at the minimum of their own objective (e.g.\ Åkerman 100\,°C: printed 117.9, global minimum \AkHundredFullGlobal{}), and random multistarts reach the global minimum in only \LandMinFrac--\LandMaxFrac\,\% of starts (Figure~\ref{fig:landscape}a). Since the choice between the full and the simplified model was made by comparing BIC values of fits that were not all at their optimum, some model-selection conclusions of the original paper (for example that the simplified Mao model has a higher BIC than the full model) should be re-read with the global optima of Table~\ref{tab:reprodyn}. For the Tao case study the printed dynamic fits could not be reproduced from the published data under any initial acid ratio; we therefore use our own re-fits of the original method as the reference throughout, which is the stronger (lower-error) baseline. The general lesson is the one that shaped our evaluation design: an in-sample $F_\mathrm{obj}$ or BIC, especially of a non-identifiable model, is not evidence that a kinetic model will predict, and kinetic parameters should be accompanied by out-of-sample validation.

\subsection{Why Pseudo-Processivity Matters, and Where Deactivation Applies}
Nishimura et al. state that the data of Bornadel et al. are affected by mass-transfer limitation and that their selectivities are therefore only apparent, but kept the irreversible model ``in the interests of simplicity''. PPK-FIP makes this choice data-driven. For the Bornadel profile the data select pseudo-processivity decisively (Step~2), the probability that a freshly formed monoester is re-acylated before release ($p_M\approx\BornPM$) is consistent with the processive fingerprint analysis of Mitchell et al. \cite{mitchell2023}, and the residual pattern it removes is one that no lumped constant can remove, because it is fixed by the composition path. Ignoring it does not merely inflate $F_\mathrm{obj}$: in the Monte-Carlo experiment with a pseudo-processive truth the two-step estimator of $S_M$ is biased by \SimSMBiasOrig\ on the log scale, whereas PPK-FIP is nearly unbiased (Step~5), and in cross-validation the processivity term alone accounts for most of the out-of-sample gain (Step~4). Conversely, when the truth is the original model, the shrinkage prior and BIC weighting keep PPK-FIP from inventing processivity (scenario S1), which is the property a generalisation of a published model must have. First-order deactivation within a batch is a profile-specific mechanism. It is selected where the late part of the profile slows down more than saturation can explain, at Åkerman 100\,°C ($k_d=\AkHundredKd$\,h$^{-1}$, well above the usual operating temperature of the immobilised lipase), and stays at the prior level ($k_d\approx10^{-3}$\,h$^{-1}$, i.e.\ no deactivation within the observed time) at 70\,°C and in the Tao and Bornadel profiles. Because the shrinkage prior switches it off where it is not needed, PPK-FIP with and without the term predict alike in cross-validation (Step~4), so the term can be carried at no cost and is available for the conditions, higher temperatures and longer or repeated runs, in which it is expected to matter.

\subsection{The Two-Step Method as a Limiting Prior}
Expressing the fingerprint step as a prior clarifies what the two-step method assumes. Fixing $S_M$ and $S_D$ is optimal only if the fingerprint estimate is exact; the time-domain data then cannot correct it even when they disagree, as they do in the Bornadel case, where the fingerprint fit has a sum of squares 2.5 times larger than the dynamic fit. A flat prior (the joint fit) uses no fingerprint information at all and over-fits small profiles, as the sparse Åkerman column of the ablation and scenario S1 of the Monte-Carlo experiment show. The calibrated prior of PPK-FIP lies between these limits, and the data confirm that the intermediate choice is best: leave-one-case-out meta-validation selected a prior three times wider than the fingerprint confidence interval ($\lambda=3$) for all six cases, and in the ablation both the tighter prior closest to the two-step method ($\lambda=1$, LOTO RMSE \AblLamOne) and the flat prior of the joint fit (\AblJoint) predict worse than PPK-FIP (\LotoAllNew). The fingerprint information is valuable and should be kept, but as information, not as a constraint. This supports the central idea of Nishimura et al. while removing its main limitation, and it also explains why the multi-task pooling of Step~4 brought no net gain: sharing the composition path across conditions is the fingerprint idea in another guise, and it pays where the conditions change only the time-scale (Tao), whereas per-profile estimation is preferable where they also change the mechanism (Åkerman).

\subsection{Mechanistic Structure versus Black-Box Learning}
Modern data-driven models are not a substitute for kinetic structure in this setting. With 7--30 sampling times per profile, the Gaussian process, MLP ensemble, pretrained TabICL foundation model and neural ODE all interpolate reasonably inside the observed range of densely sampled profiles, and the neural ODE interpolates the Mao batches more closely by following the batch-1 plateau that mass conservation excludes. All of them fail in extrapolation (Step~6), because nothing in them encodes mass conservation, the Ping Pong rate law or the sequential nature of the esterifications; the neural ODE's prediction that the monoester re-appears late in the Bornadel run (Figure~\ref{fig:fcst}) illustrates what is lost. The improvements of PPK-FIP come instead from the mechanism the data actually require and from better use of the fingerprint information, and the exact time reparameterisation is what makes global optimisation, hundreds of cross-validation refits and Monte-Carlo studies cheap enough to do routinely. Machine learning is most useful \emph{on top of} the mechanistic model: a Gaussian-process discrepancy absorbs the systematic and digitisation errors that no kinetic model should try to explain, giving the lowest interpolation error of all methods (PPK-FIP-GP), while the mechanistic mean guarantees sensible behaviour outside the data. That the same discrepancy term also rescues the interpolation of the original two-step model, but not its extrapolation, is the clearest demonstration in the paper that interpolation accuracy and mechanistic adequacy are different things.

\subsection{Implications for Biolubricant Process Design}
For the design of processes that convert biomass-derived fatty acids and TMP into biolubricants, the results have three practical consequences. Apparent selectivities obtained in a mass-transfer-limited reactor describe that reactor, not the enzyme: the Bornadel values ($S_M=0.909$) would predict a different monoester-to-diester balance from the processivity-corrected values ($S_M=\BornSM$) if transferred to a reactor with better mixing or smaller particles, so a processivity correction should accompany any such transfer. Within a system whose conditions change only the time-scale, the composition path can be transferred between conditions (Tao water activities), which reduces the experimental effort needed to characterise a new condition to the estimation of the rate constants. And kinetic parameters intended for design should be reported together with an out-of-sample check such as leave-one-time-point-out cross-validation, which here cost seconds per fold with the exact reparameterisation and which separates models that describe from models that predict.

\subsection{Scope of the Results}
The results should be read within the following scope. All profiles were digitised from figures of the source publications, so the residual level includes digitisation error and differs between studies, and the Mao batch-1 plateau is treated as a measurement feature rather than as kinetics; this is the kind of error that the Gaussian-process discrepancy of PPK-FIP-GP is designed to absorb. The statistical comparisons rest on leave-one-time-point-out validation with \NLotoFolds{} folds and \NFitPoints{} observations; the forward-extrapolation protocol, with one split per case, characterises behaviour beyond the data window and the information content of half-profiles rather than ranking methods. The lumped constants remain empirical, as the original authors emphasised, and the full set is determined only up to the ratios of Figure~\ref{fig:landscape}b, so interpretation should focus on the predictions and on the identifiable parameters of the composition path ($S_M$, $S_D$, $p_M$, $p_D$) and of deactivation ($k_d$). Pseudo-processivity is an effective description of mass-transfer effects; its parameters belong to the reactor and agitation conditions and should be re-estimated when these change, which the fingerprint prior and the exact reparameterisation make inexpensive. Finally, the prior width was meta-validated on the six case studies of this reaction class; the same procedure transfers directly to a new class of lipase-catalysed process, where $\lambda$ is selected from that class's own cases.
